\section{Applying fracture to the height-three overlap}\label{sec:application}

The coefficient arguments have now produced the maps required by the
finite construction of \Cref{prop:fracture}.  We apply it with one
choice of local maps and compatibility homotopy, obtaining both the
filtration and its attachment identification.  The local basis results
used here were proved without the rational-obstruction theorem, which
enters only in the final subsection.

\subsection{Checking the fracture data}

\begin{proof}[Proof of \Cref{thm:main} and the height-three instance of
\Cref{thm:attachment-identification}]
Put \(T=L_{K(3)}S\), \(X=L_2T\), and \(R_1=L_{K(1)}S\).
Recall \(Q=L_0S\simeq H\Qp\), and set
\[
 W=W_1\vee W_2=L_2S\vee\Sigma^{-1}L_2S.
\]
\Cref{h3:thm:height-two-test} supplies the actual unit and determinant
map \(w=(1,\zeta):W\to X\), a \(K(2)\)-equivalence.  This verifies
\textup{(i)} of \Cref{prop:fracture}.  Put \(Y=\cofib(w)\) and
write \(q:X\to Y\).

For the middle block, set
\[
 U_{\mathrm{mid}}=U_3\vee U_4=\Sigma^{-3}L_1S\vee\Sigma^{-4}L_1S.
\]
\Cref{h3:thm:height-one-maps} supplies the equivalence of actual
\(R_1\)-module maps
\begin{equation}\label{eq:height-one-basis}
 (1,\zeta,\delta,\zeta\delta):
 R_1\vee\Sigma^{-1}R_1\vee\Sigma^{-3}R_1\vee\Sigma^{-4}R_1
 \xrightarrow{\simeq}L_{K(1)}T.
\end{equation}
Its first two components are the localizations of the same maps from
\(S\) used in \(w\).  Their restrictions are the same unit and
determinant maps, and the module localization adjunction
\eqref{eq:module-localization} identifies their extensions to \(R_1\),
including the required homotopies.  The fourth component is the specified
product \(\zeta\delta\).  Localizing the cofiber sequence at \(K(1)\)
therefore identifies
\(L_{K(1)}Y\) with the cofiber of the first two components of
\eqref{eq:height-one-basis} using those specified maps.
Consequently the last two components of
\eqref{eq:height-one-basis}, followed by the quotient, give
the required equivalence
\begin{equation}\label{eq:kappa-application}
 \kappa:L_{K(1)}U_{\mathrm{mid}}\xrightarrow{\simeq}L_{K(1)}Y.
\end{equation}
They are represented by the projections of \(\delta\) and
\(\zeta\delta\).

For the rational part of the middle map,
\Cref{h3:thm:rational-compatibility} gives
\[
 \pi_*L_0T=\Lambda_{\Qp}(\zeta,\alpha_3,\alpha_5),
 \qquad |\zeta|=-1,\quad|\alpha_3|=-3,\quad|\alpha_5|=-5,
\]
and identifies the actual rationalized localization map by
\begin{equation}\label{eq:rational-images}
 \alpha_3\longmapsto c\delta,\qquad
 \zeta\alpha_3\longmapsto c\zeta\delta
 \quad\text{for one }c\in\Qp^\times.
\end{equation}
Set
\[
 \beta_3=c^{-1}\alpha_3.
\]
This normalization is made only in the rational source.
No condition \(c\in\Zp^\times\) is required.
The projections of \(\beta_3\) and \(\zeta\beta_3\) to
\(L_0Y\) define a \(Q\)-module map
\[
 r:L_0U_{\mathrm{mid}}=\Sigma^{-3}Q\vee\Sigma^{-4}Q\longrightarrow L_0Y.
\]
By \eqref{eq:rational-images}, their images in
\(L_0L_{K(1)}Y\) are the projections of
\(\delta\) and \(\zeta\delta\).  These are also the images,
under \(L_0\kappa\), of the two unit generators of \(L_0U_{\mathrm{mid}}\).
The maps in \eqref{eq:compatibility} therefore agree on the
generators of this finite free \(Q\)-module.  An identifying
homotopy of \(Q\)-module maps exists; choose one as the
compatibility datum in \textup{(ii)} of \Cref{prop:fracture}.

It remains to check the final quotient.  Since \(L_0L_iS=Q\)
for \(i=1,2\), rationalization of \(w\) is the map from
\(Q\vee\Sigma^{-1}Q\) represented by \(1,\zeta\).
These are two distinct members of the exterior basis in
\Cref{h3:thm:rational-compatibility}, so this map is injective
on rational homotopy.  Its cofiber \(L_0Y\) has one copy of
\(\Qp\) in each degree
\[
 -3,\ -4,\ -5,\ -6,\ -8,\ -9.
\]
There are no suspended kernel terms.  The map \(r\) is likewise
injective on homotopy and includes the first two lines, represented
by \(\beta_3,\zeta\beta_3\).

Put
\[
\begin{aligned}
 V&=V_5\vee V_6\vee V_7\vee V_8\\
  &=\Sigma^{-5}Q\vee\Sigma^{-6}Q
       \vee\Sigma^{-8}Q\vee\Sigma^{-9}Q.
\end{aligned}
\]
Define \(b:V\to L_0X\simeq L_0T\) by the four ordered classes
\begin{equation}\label{eq:rational-basis-application}
 \alpha_5,\qquad \zeta\alpha_5,\qquad
 \beta_3\alpha_5,\qquad\zeta\beta_3\alpha_5.
\end{equation}
They represent the remaining quotient lines in degrees
\(-5,-6,-8,-9\).  The composite \eqref{eq:quotient-basis}
is therefore an isomorphism on all homotopy groups, and hence
an equivalence of rational modules.  This verifies
\textup{(iii)}.

Carry out \Cref{prop:fracture} with this single choice of
\(w,\kappa,r,b\) and compatibility homotopy, choosing coherent
inverse data for \(L_{K(2)}w\) and for \(L_Ja\) at the steps where
those equivalences are established.  This gives the same
\(h,F_i,a,Z,t,e\) used by \Cref{thm:attachment-identification},
with exactly the ordered cofibers in \Cref{thm:main}.
The terminal projection \(e:F_8\to X\) is an equivalence.  Its
restriction to \(F_1=L_2S\) is the first component of \(w\),
namely the canonical map \(L_2(S\to T)\), proving \Cref{thm:main}.
For these same choices, \Cref{thm:attachment-identification}
identifies all eight connecting maps by its two signed roofs.
This proves its height-three instance without replacing the
filtration or choosing new local maps.
\end{proof}

\begin{remark}
The exact compatibility in \eqref{eq:rational-images} is needed
for this application.  Rational dimensions alone would not give
the homotopy in \eqref{eq:compatibility}; independently chosen
degree-three and degree-four maps would not ensure that the same
normalization works for the product with \(\zeta\).
\end{remark}
